\section{Implementation Surface}
\label{sec:sources-wasm4pm-implementation}

\subsection{Source Files}
\label{sec:sources-wasm4pm-impl-source-files}

The crate root is \texttt{src/lib.rs}, declaring 17 public modules. The conformance module
(\texttt{src/conformance.rs}, 1,229 lines) is the largest single file and implements: the
\texttt{ConformanceVerdict} enum with \texttt{FullyConforming} / \texttt{PartiallyConforming} /
\texttt{NonConforming} variants; \texttt{ConformanceVerdicts} aggregate container with
\texttt{aggregate\_fitness}, \texttt{aggregate\_precision}, \texttt{admitted\_cases}, and
\texttt{total\_cases}; \texttt{TokenReplayEngine} with source/sink token accounting and the van
der Aalst fitness equation $f = 0.5(1 - m/c) + 0.5(1 - r/p)$; \texttt{AlignmentEngine} with
real A* search using \texttt{BinaryHeap}, a 5000-iteration cap, and compile-time metric
validation via \texttt{Between01<NUM,DEN>} and \texttt{Metric<KIND,NUM,DEN>} const generics; and
\texttt{RuntimeBetween01} with NaN/Inf/bounds guards for runtime \texttt{f64} values.

The mining module (\texttt{src/mining/mod.rs}, 847 lines, sealed receipt
\texttt{MINING\_RENDER\_RECEIPT.md}) was manufactured by the ggen manufacturing machinery and
implements: \texttt{ProcessModel} union type; three witness structs (\texttt{AlphaWitness},
\texttt{InductiveWitness}, \texttt{HeuristicsWitness}) each implementing \texttt{Lattice<T>};
\texttt{Admitted} state enum; and four public API functions returning
\texttt{Evidence<ProcessModel,~Admitted,~W>} containers. The crypto module
(\texttt{src/crypto.rs}) self-implements BLAKE3 with Merkle-tree chaining, SHA-256, SHA-512,
ChaCha20, Twisted-Edwards Curve25519 field arithmetic (\texttt{CurvePoint},
\texttt{FieldElement}), Ed25519 signature verification (\texttt{verify\_ed25519\_signature}, RFC
8032 compliant with cofactor-cleared equation check), and \texttt{verify\_jcs\_receipt\_signature}
(JCS RFC 8785 canonicalization + Ed25519 verification). The graduation module
(\texttt{src/graduation.rs}) implements the \texttt{GraduateToWasm4pm} trait, bridging
\texttt{wasm4pm-compat} \texttt{GraduationCandidate} types into the execution engine.

Additional modules: \texttt{src/evidence.rs} (Evidence carrier and lattice infrastructure),
\texttt{src/ltl.rs} (LTL/Declare constraint evaluation: \texttt{DeclareRule} with
\texttt{Precedence}/\texttt{Response} variants), \texttt{src/sandbox.rs} (\texttt{GasMeter},
\texttt{RecursionGuard}, \texttt{execute\_oblivion\_protocol}), \texttt{src/otel.rs}
(\texttt{OtelTrace} / \texttt{OtelSpan} with \texttt{blake3\_receipt}, \texttt{witness\_id},
\texttt{token\_state} lifecycle metadata), \texttt{src/ocel\_v2.rs}
(\texttt{ZeroCopyOcelV2} zero-allocation binary parser), and \texttt{src/lifecycle/mod.rs}
(\texttt{ProcessIntelligence<S>} 13-state MAPE-K typestate machine with all five transition proof
witness types).

\subsection{Commands}
\label{sec:sources-wasm4pm-impl-commands}

The crate builds with \texttt{cargo build -{}-release} and all 31 tests pass via \texttt{cargo
test} in 10.60 seconds. The mining module was manufactured and compiled in isolation in 1.73
seconds (\texttt{cargo build -{}-release}, wasm4pm v30.1.2). No other build tooling is required
beyond the Rust toolchain and the \texttt{wasm4pm-compat} path dependency. The research verdict
specifies a WASM runtime spine with a 4~GB linear memory layout; deployment to a WASM host
requires the \texttt{cdylib} artifact linked into a WASM runtime (host-side integration is not
part of this crate).

\subsection{Data and Ontology Surfaces}
\label{sec:sources-wasm4pm-impl-data-ontology}

The \texttt{mining-authority-ontology.md} documents the evidence type registry: four evidence
types (Inductive, Heuristics, Alpha, DFG), all lattice-compliant. The
\texttt{OBJECT\_CENTRIC\_RUNTIME.md} documents the object-centric runtime design. The
\texttt{ZeroCopyOcelV2} parser in \texttt{src/ocel\_v2.rs} enforces OCEL~2.0 binary format
compliance with magic/version/offset validation. \texttt{OtelTrace} and \texttt{OtelSpan} in
\texttt{src/otel.rs} implement the OpenTelemetry span model augmented with process-intelligence
fields (\texttt{blake3\_receipt}, \texttt{witness\_id}, \texttt{token\_state}).

\subsection{Templates and Rendered Artifacts}
\label{sec:sources-wasm4pm-impl-templates}

The mining module (\texttt{src/mining/mod.rs}) was materialized by the ggen manufacturing
machinery from templates sourced at
\texttt{sources/wasm4pm-compat/compat/templates/mining/}: \texttt{alpha\_miner.rs.j2},
\texttt{inductive\_miner.rs.j2}, and \texttt{heuristics\_miner.rs.j2}, with bound template
variables \texttt{module\_name}~=~``mining'', \texttt{witness\_markers}~=~[``VanDerAalst1989'',
``Leemans2013'', ``Weijters2011''], \texttt{evidence\_carriers}~=~[``ProcessModel'',
``Admitted''], and \texttt{lifecycle\_states}~=~[``Initial'', ``Discovered'', ``Sealed''].
The sealed receipt (\texttt{MINING\_RENDER\_RECEIPT.md}, status SEALED,
2026-06-01T00:00:00Z) documents the complete manufacturing chain from specification extraction
through template application, artifact materialization, zero-warning compilation verification,
and authority sealing. This receipt is the primary evidence artifact demonstrating that the mining
module was not hand-coded but manufactured from an authoritative specification.
